\section{Finite realization of instruments with interior posterior states}
\label{sec:r74-instruments}

The preceding response problems have a terminal observation. We now allow
an outcome at every call, state-dependent outcome probabilities, and subsequent
commands selected from the observed history. The sufficient condition is on
each posterior state, not merely on the channel obtained by forgetting the
outcome. The construction uses the classical positive-realization principle
of an invariant finitely generated cone \cite{Heller65,MonrasWinter14}.
For quantum instruments, preservation of a simulation polytope and its
composition under adaptive operations are developed explicitly by
Okay~\cite{Okay26}. We use that established mechanism and prove a uniform
posterior condition which supplies one finite physical dictionary, an explicit
dimension and noise bound, and a counted pure initialization. The accompanying
counterexample distinguishes this condition from complete mixing of the
unconditioned channel.

Let $\mathcal D_h$ be the density matrices on $\mathbb C^h$, $h\ge2$.
For $0<\eta\le1$ put
\[
 \mathcal D_{h,\eta}
   =(1-\eta)\mathcal D_h+\eta I/h
   =\{\rho\in\mathcal D_h:\rho\succeq\eta I/h\}.
\]
The following elementary geometric estimate is included in full. Nets of
pure states and inner approximations of the density-matrix body are classical
tools; see Aubrun--Szarek~\cite{AubrunSzarek17}. In particular the bound below
is not asserted optimal in the dimension $h$.

\begin{lemma}[A finite physical dictionary]\label{lem:r74-state-dictionary}
There are rank-one density matrices $\sigma_1,\ldots,\sigma_K$ such that
\begin{equation}\label{eq:r74-state-dictionary}
 \mathcal D_{h,\eta}\subseteq
 P:=\operatorname{conv}\{\sigma_1,\ldots,\sigma_K\}
 \subseteq\mathcal D_h,
 \qquad K\le(4h/\eta)^{h-1}.
\end{equation}
The same dictionary works for every map whose normalized positive outputs
belong to $\mathcal D_{h,\eta}$. If $\eta$ is real algebraic, the dictionary
can be chosen algebraic and constructed by a finite real-algebraic procedure.
\end{lemma}
\begin{proof}
For rank-one projectors use the metric
\[
 d(\psi,\varphi)=\sqrt{1-|\langle\psi,\varphi\rangle|^2}
   =2^{-1/2}\bigl\||\psi\rangle\langle\psi|
                   -|\varphi\rangle\langle\varphi|\bigr\|_2.
\]
It is a metric on projective space, since the last expression is the
Hilbert--Schmidt distance up to a fixed factor. The invariant probability
measure of a ball of radius $r\in[0,1]$ is $r^{2h-2}$. To verify the formula,
fix its center to the first coordinate vector and normalize $h$ independent
standard complex Gaussian variables. Their squared absolute values are
independent exponentials. After division by their sum the resulting weights
are uniform on the probability simplex, with density $(h-1)!$. The region
where the first weight is at least $1-r^2$ is a simplex scaled by $r^2$
in its $h-1$ free coordinates, giving the stated measure. This argument also
derives the invariant law, since the joint Gaussian density is unitary
invariant.

Take a maximal family with mutual distances greater than
$r=\sqrt{\eta/h}$. The balls of radius $r/2$ are disjoint up to their
boundaries, so the family is finite and has at most
$(2/r)^{2h-2}=(4h/\eta)^{h-1}$ points. Maximality gives a covering by
balls of radius $r$. Let the corresponding projectors be $\sigma_i$.
For any Hermitian $H$, choose a unit eigenvector $\psi$ for its largest
eigenvalue $a$, and let $z$ be its smallest eigenvalue. A covering point
$\varphi$ satisfies $|\langle\psi,\varphi\rangle|^2\ge1-\eta/h$.
Consequently
\[
 \max_i\tr(H\sigma_i)\ge a-(\eta/h)(a-z)
 \ge(1-\eta)a+(\eta/h)\tr H.
\]
The second inequality follows from $ha-\tr H\ge a-z$.
The right side is the support function of $\mathcal D_{h,\eta}$.
Separation of compact convex sets proves the first containment in
\eqref{eq:r74-state-dictionary}; the second follows because every
$\sigma_i$ is a density matrix.

For algebraic $\eta$, choose the separated projectors successively with
algebraic real and imaginary entries. At each step the existence of a
rank-one projector at distance greater than $r$ from all selected ones
is a finite system of real polynomial conditions over the real algebraic
numbers. Quantifier elimination decides it and, when it is feasible,
returns an algebraic witness. The same packing bound forces termination;
at termination the selected projectors cover the whole projective space.
Thus the cardinality and containment bounds are unchanged. This finite
construction has no efficiency assertion.
\end{proof}

For each command $a$ in a finite set $A$, let
$\{\Phi_{a,y}:y\in Y_a\}$ be a finite quantum instrument: each map is
completely positive and their sum is trace preserving. Its primitives and
initial density matrix $\rho_0$ are fixed and known. Assume
\begin{equation}\label{eq:r74-posterior-positivity}
 \Phi_{a,y}(\rho)\succeq
       (\eta/h)\tr(\Phi_{a,y}(\rho))I
 \quad(\rho\in\mathcal D_h,\ a\in A,\ y\in Y_a).
\end{equation}
Zero-probability branches are allowed. For a positive branch this is precisely
the assertion that its normalized posterior belongs to
$\mathcal D_{h,\eta}$.

A classical instrument simulator has nonnegative matrices $M_{a,y}$ on one
fixed finite label set, with $\sum_yM_{a,y}$ row-stochastic for every $a$.
The entry $M_{a,y}(i,j)$ is the joint probability of observing $y$ and retaining
label $j$ from label $i$. Its initial label is prescribed pure. Its internal
label is not revealed to a policy which originally observed only instrument
outcomes.

\begin{theorem}[Uniform posterior positivity gives finite exact memory]
\label{thm:r74-instrument}
Under \eqref{eq:r74-posterior-positivity}, there is a single classical
instrument simulator, with the prescribed pure start, using at most
\begin{equation}\label{eq:r74-instrument-bound}
 1+\left\lfloor(4h/\eta)^{h-1}\right\rfloor
\end{equation}
live labels. It reproduces the joint law of every finite command--outcome
history and every subsequent terminal POVM. The same simulator works under
every policy that selects commands and a stopping time from the visible
history, including randomized policies. It therefore preserves every
almost surely finite stopped history law, every bounded score on that law,
and the extended expectation of every nonnegative history cost.

The bound does not depend on the number of commands or outcomes. The sizes
of the transition and terminal tables do depend on those numbers. For
$\eta=1$ at most two labels suffice; for $h=1$ one label suffices.
\end{theorem}
\begin{proof}
Choose the dictionary from Lemma~\ref{lem:r74-state-dictionary} and adjoin
one label $0$, represented by $\sigma_0=\rho_0$. For each label $i$,
command $a$ and outcome $y$ put $p_{i,a,y}=\tr\Phi_{a,y}(\sigma_i)$.
If this is positive, the normalized output lies in $P$, so choose
nonnegative numbers $v_{i,a,y,j}$, $1\le j\le K$, of sum one, with
\[
 \Phi_{a,y}(\sigma_i)
     =p_{i,a,y}\sum_{j=1}^K v_{i,a,y,j}\sigma_j.
\]
Set $M_{a,y}(i,j)=p_{i,a,y}v_{i,a,y,j}$ for $j\ge1$ and
$M_{a,y}(i,0)=0$. On a zero-probability branch set the entire row to zero;
its positive output operator is zero because its trace is zero. The sum
over $y,j$ of each row is one, by trace preservation. Thus these are valid
classical instrument tables and the initializer is one of the counted labels.
There is no uncounted random preparation or input-reading step.

Start at label $0$. Induction gives the operator identity
\begin{equation}\label{eq:r74-instrument-intertwining}
 \Phi_{a_n,y_n}\cdots\Phi_{a_1,y_1}(\rho_0)
   =\sum_{j=0}^K
       (e_0^{\mathsf T}M_{a_1,y_1}\cdots M_{a_n,y_n})_j\sigma_j.
\end{equation}
For the empty history this is $\rho_0=\sigma_0$. For the inductive step
apply the next linear map and then the defining convex expansions.
Taking traces proves every joint history probability. Pairing with any
terminal effect $E_c$ proves the terminal law with decoder
$D(i,c)=\tr(E_c\sigma_i)$, which is row-stochastic for a POVM.
It also covers a history-dependent choice of terminal POVM.

For an adaptive policy, multiply each history probability by the policy's
conditional command and stopping probabilities. They are the same functions
of the visible history in both experiments, so the equality holds on every
finite cylinder and on every finite stopped history. These cylinders generate
the path sigma-field; countable additivity gives equality of the full path
laws. An almost surely finite stopped law is their countable union. Equality
of bounded score expectations and, by monotone convergence, nonnegative cost
expectations follows. A policy's own persistent memory is additional to the
simulator labels; the construction does not grant it observation of those
labels.

When $\eta=1$, every positive posterior is $I/h$, so one recurrent label
represented by $I/h$ and the initializer suffice. When $h=1$, the state is
fixed and the instrument consists of its classical outcome probabilities.
\end{proof}

The theorem compiles a fixed known preparation and set of primitives. It
does not implement a measurement which converts an arbitrary unknown quantum
input into a classical label, and it does not give shared parameter-independent
tables for an unknown model family. Those are different quantifiers.

\begin{corollary}[Noise applied after every observed branch]
\label{cor:r74-noisy-instruments}
Let $\{\Psi_{a,y}\}$ be arbitrary finite quantum instruments on
$\mathbb C^h$. For any $0<\eta\le1$, the instruments
\begin{equation}\label{eq:r74-branch-noise}
 \Phi_{a,y}(X)=(1-\eta)\Psi_{a,y}(X)
            +\eta\tr(\Psi_{a,y}(X))I/h
\end{equation}
have the finite exact simulator of Theorem~\ref{thm:r74-instrument}.
Outcome probabilities can depend on the incoming state, and the commands
need not commute or share a stationary state.
\end{corollary}
\begin{proof}
The two summands are completely positive; their traces sum to the original
branch trace. The command sums remain trace preserving. Positivity of the
first summand gives \eqref{eq:r74-posterior-positivity} directly.
\end{proof}

Once a dictionary is chosen, its transition coefficients are solutions of
finite linear feasibility problems. If the primitives, initial state and
$\eta$ are algebraic, the lemma supplies an algebraic dictionary and
linear feasibility gives algebraic tables. The required terminal decoders
are algebraic when the supplied terminal POVMs are algebraic. This observation is not a claim
of efficient net construction, small program size, or a matching lower bound
for the unrestricted stochastic memory of each individual experiment.
Proposition~\ref{prop:r74-dictionary-lower} records the matching small-noise
power for the narrower problem of one physical dictionary covering all of
$\mathcal D_{h,\eta}$. Stronger dimension estimates for such convex
approximations are already known \cite{AubrunSzarek17}; extension complexity
is another distinct resource \cite{Fawzi21}.
